\section{Bài tập 7: Full-wave Rectifier}

Cho mạch chỉnh lưu cầu toàn sóng sử dụng máy biến áp như hình bên dưới. Nguồn điện áp đầu vào là nguồn sin xoay chiều $V_1$ có các thông số:
\begin{itemize}[leftmargin=1.5em]
  \item \texttt{VOFF} = $0\text{ V}$, \texttt{VAMPL} = $220\text{ V}$, \texttt{FREQ} = $50\text{ Hz}$, \texttt{AC} = $1$.
\end{itemize}
Máy biến áp $TX_1$ có tỷ số biến áp $N_1/N_2 = 10$. Cầu chỉnh lưu gồm 4 diode $D_1, D_2, D_3, D_4$ (\texttt{Dbreak}) cấp nguồn cho điện trở tải $R_1 = 1\text{ k}\Omega$.

Giả thiết các diode có điện áp phân cực thuận ngưỡng là $V_F = 0.7\text{ V}$. Hãy viết biểu thức hiệu điện thế $V_{AB}(t)$ tại cuộn thứ cấp của biến áp và $V_{CD}(t)$ trên điện trở tải. Sau đó thực hiện phân tích miền thời gian (Transient) trên PSpice.

\given{ex7_given}{0.65\textwidth}{Mạch chỉnh lưu toàn sóng dạng cầu có máy biến áp (Hình 2.23)}

\subsection{Tính toán lý thuyết}
Điện áp tức thời ở cuộn sơ cấp máy biến áp:
\begin{equation}
  v_1(t) = 220 \sin(2\pi \cdot 50 t) = 220 \sin(100\pi t)\text{ (V)}
\end{equation}

\noindent Tỷ số máy biến áp là $N_1/N_2 = 10$, do đó biên độ điện áp xoay chiều tại cuộn thứ cấp giữa hai điểm A và B là:
\begin{equation}
  V_{AB,\max} = \frac{V_m}{N_1/N_2} = \frac{220\text{ V}}{10} = 22\text{ V}
\end{equation}
\noindent Phương trình hiệu điện thế $V_{AB}(t)$:
\begin{equation}
  V_{AB}(t) = 22 \sin(100\pi t)\text{ (V)}
\end{equation}

\noindent Trong mạch chỉnh lưu cầu, ở mỗi nửa chu kỳ luôn có 2 diode mắc nối tiếp nhau cùng dẫn điện:
\begin{itemize}[leftmargin=1.5em]
  \item Bán kỳ $V_{AB} > 0$ (A dương hơn B): cặp $D_2$ và $D_3$ dẫn thuận ($D_2$: anode A, cathode C; $D_3$: anode D, cathode B). Dòng điện đi từ nút A qua $D_2 \to$ nút C $\to$ tải $R_1$ (từ C xuống D) $\to$ nút D $\to D_3 \to$ nút B.
  \item Bán kỳ $V_{AB} < 0$ (B dương hơn A): cặp $D_1$ và $D_4$ dẫn thuận ($D_1$: anode B, cathode C; $D_4$: anode D, cathode A). Dòng điện đi từ nút B qua $D_1 \to$ nút C $\to$ tải $R_1$ (từ C xuống D) $\to$ nút D $\to D_4 \to$ nút A.
\end{itemize}
Tổng sụt áp trên hai diode dẫn là:
\begin{equation}
  2V_F = 2 \times 0.7\text{ V} = 1.4\text{ V}
\end{equation}

\noindent Do đó, dòng điện qua tải luôn chảy theo cùng một chiều (từ C sang D) trong cả hai nửa chu kỳ, và biên độ đỉnh của điện áp ngõ ra $V_{CD}$ là:
\begin{equation}
  V_{CD,\max} = V_{AB,\max} - 2V_F = 22\text{ V} - 1.4\text{ V} = 20.6\text{ V}
\end{equation}
\noindent Phương trình hiệu điện thế ngõ ra $V_{CD}(t)$:
\begin{equation}
  V_{CD}(t) = \max\left(0,\, 22|\sin(100\pi t)| - 1.4\right)\text{ (V)}
\end{equation}

\begin{dapsobox}
  \begin{itemize}[leftmargin=1.5em]
    \item $V_{AB}(t) = 22 \sin(100\pi t)\text{ (V)}$
    \item $V_{CD}(t) = \max\left(0,\, 22|\sin(100\pi t)| - 1.4\right)\text{ (V)}$
  \end{itemize}
\end{dapsobox}

\subsection{Cấu hình và Mô phỏng PSpice}
\subsubsection*{1. Tính toán thông số mô phỏng:}
\begin{itemize}[leftmargin=1.5em]
  \item Chu kỳ của dạng sóng xoay chiều $V_{AB}$:
  \begin{equation}
    T = \frac{1}{f} = \frac{1}{50\text{ Hz}} = 0.02\text{ s} = 20\text{ ms}
  \end{equation}
  \item Thời gian cần thiết để mô phỏng $10$ chu kỳ của dạng sóng $V_{AB}$ (\textit{Run to time}):
  \begin{equation}
    T_{\text{sim}} = 10 \times T = 10 \times 20\text{ ms} = 200\text{ ms}
  \end{equation}
  \item Để tần số lấy mẫu cao gấp 10 lần tần số tín hiệu ($f_s = 10 \times 50\text{ Hz} = 500\text{ Hz}$), khoảng thời gian giữa hai điểm lấy mẫu liên tiếp (\textit{Maximum step size}) là:
  \begin{equation}
    \Delta t = \frac{1}{f_s} = \frac{1}{500\text{ Hz}} = 2\text{ ms} \quad (\text{để đường cong mịn hơn có thể đặt } 0.1\text{ ms})
  \end{equation}
\end{itemize}

\subsubsection*{2. Tiến hành mô phỏng:}
\begin{itemize}[leftmargin=1.5em]
  \item Đặt biến áp: vào \texttt{Place > PSpice Component > Modeling Application... > System Modules > Transformer}, thiết lập tỷ số biến áp $10:1$.
  \item Đặt đầu dò vi sai (\textbf{Voltage Differential Marker}): cặp que đo vi sai thứ nhất tại hai đầu A và B của cuộn thứ cấp; cặp que đo vi sai thứ hai tại hai đầu C và D của điện trở tải $R_1$.
  \item Chạy mô phỏng Transient với \textit{Run to time} = $100\text{ ms}$ (tương ứng 5 chu kỳ của $V_{AB}$) và \textit{Maximum step size} = $0.1\text{ ms}$.
\end{itemize}

\simresult
\simfig[8cm]{ex7_sim1}{Đồ thị so sánh dạng sóng ngõ vào $V_{AB}$ và ngõ ra $V_{CD}$ trên cùng một cửa sổ (con trỏ 1 tại $t = 5\text{ ms}$, con trỏ 2 tại $t = 15\text{ ms}$)}

\begin{table}[H]
  \centering
  \caption{Số liệu đọc bằng con trỏ để kiểm chứng phương trình $V_{AB}(t)$ và $V_{CD}(t)$}
  \begin{tabular}{lcccccc}
    \toprule
    & \multicolumn{3}{c}{\textbf{Lý thuyết}} & \multicolumn{3}{c}{\textbf{PSpice (cursor)}} \\
    \cmidrule(lr){2-4} \cmidrule(lr){5-7}
    \textbf{Vị trí đo} & $t\text{ (ms)}$ & $V_{AB}\text{ (V)}$ & $V_{CD}\text{ (V)}$ & $t\text{ (ms)}$ & $V_{AB}\text{ (V)}$ & $V_{CD}\text{ (V)}$ \\
    \midrule
    Đỉnh dương của $V_{AB}$ & $5$  & $22.0$  & $20.6$ & $5.0$ & $21.317$ & $19.848$ \\
    Điểm qua $0$            & $10$ & $0$     & $0$    & $10$ & $\approx 0$ & $\approx 0$ \\
    Đỉnh âm của $V_{AB}$    & $15$ & $-22.0$ & $20.6$ & $15.0$ & $-22.585$ & $21.113$ \\
    \bottomrule
  \end{tabular}
\end{table}

\subsection{So sánh và nhận xét}
\begin{table}[H]
  \centering
  \caption{So sánh các thông số mạch chỉnh lưu toàn sóng giữa Lý thuyết và PSpice}
  \begin{tabular}{lccc}
    \toprule
    \textbf{Đại lượng} & \textbf{Lý thuyết (Theory)} & \textbf{Mô phỏng (PSpice)} & \textbf{Sai lệch} \\
    \midrule
    Biên độ cực đại $V_{AB,\max}\text{ (V)}$ & $22.0$ & $22.585$ & $2.7\%$ \\
    Biên độ cực đại $V_{CD,\max}\text{ (V)}$ & $20.6$ & $21.113$ & $2.5\%$ \\
    Tần số gợn sóng $f_{\text{out}}\text{ (Hz)}$ & $100.0$ & $100.0$ & $0\%$ \\
    Chu kỳ ngõ ra $T_{\text{out}}\text{ (ms)}$   & $10.0$ & $10.0$ & $0\%$ \\
    \bottomrule
  \end{tabular}
\end{table}

\noindent\textbf{Nhận xét:}
\begin{itemize}[leftmargin=1.5em]
  \item Mạch chỉnh lưu toàn sóng tận dụng được cả hai nửa chu kỳ của tín hiệu xoay chiều, do đó tần số gợn sóng đầu ra tăng gấp đôi ($2f = 100\text{ Hz}$) và chu kỳ giảm còn một nửa ($10\text{ ms}$) so với mạch nửa sóng.
  \item Điện áp đỉnh ngõ ra thấp hơn điện áp đỉnh thứ cấp một lượng đúng bằng $2V_F \approx 1.4\text{ V}$ do luôn có hai diode nối tiếp cùng dẫn dòng.
  \item Hiệu suất biến đổi năng lượng của mạch cầu cao hơn nhiều so với mạch chỉnh lưu nửa sóng và việc lọc phẳng điện áp ở giai đoạn sau sẽ dễ dàng hơn.
\end{itemize}
